\begin{onehalfspace}

La scelta degli strumenti di simulazione non è neutrale. In informatica quantistica, il framework adottato determina quali modelli di rumore sono disponibili, verso quale hardware è possibile traslare i circuiti e con quali risultati della letteratura è possibile confrontarsi direttamente. Questo capitolo descrive gli strumenti utilizzati nella parte sperimentale di questa tesi, motivando ciascuna scelta e collocandola rispetto alle alternative disponibili.

% -----------------------------------------------
\subsection{Qiskit: il Framework Quantistico di Riferimento}
\label{sec:qiskit}

\textbf{Qiskit} (\textit{Quantum Information Science Kit}) è un \ac{SDK} open-source sviluppato da \ac{IBM} Quantum per la programmazione, la simulazione e l'esecuzione di algoritmi quantistici \cite{qiskit2023}. Rilasciato pubblicamente nel 2017, offre integrazione diretta con l'hardware \ac{IBM} ed è uno dei framework ampiamente usati nella ricerca su circuiti quantistici a corto termine \cite{skosana2021, shor_hpec2023}.

\subsubsection{Architettura del Framework}

Qiskit è organizzato in tre componenti principali, ciascuno con un ruolo distinto nel ciclo di vita di un esperimento quantistico:

\begin{itemize}

    \item \textbf{Qiskit (core)} — il nucleo del framework, già noto come Qiskit Terra. Fornisce le primitive per costruire i circuiti, compilarli verso un backend target tramite \texttt{transpile} e accedere alle interfacce \texttt{Sampler} ed \texttt{Estimator}. Gestisce la pipeline dalla descrizione astratta fino all'invio al simulatore o all'hardware reale.

    \item \textbf{Qiskit Aer} — il modulo di simulazione classica ad alte prestazioni. Implementa diversi metodi di simulazione esatta e approssimata, con supporto nativo per modelli di rumore realistici. È il componente centrale di questa tesi e viene descritto in dettaglio nella Sezione~\ref{sec:qiskit_aer}.

    \item \textbf{Qiskit IBM Runtime} — l'interfaccia per l'esecuzione su processori \ac{IBM} Quantum reali tramite cloud (\texttt{qiskit-ibm-runtime}). Fornisce le primitive \texttt{Sampler} e \texttt{Estimator} ottimizzate per l'hardware reale, con supporto per tecniche di soppressione del rumore come il \textit{Dynamical Decoupling} e il \textit{Gate Twirling} \cite{ibm_shor_tutorial}.

\end{itemize}

\subsubsection{Perché Qiskit e non un'alternativa}
\label{sec:confronto_framework}

La scelta di Qiskit è stata formalizzata nell'incontro con il relatore del 18 marzo 2026, in seguito all'analisi del tutorial ufficiale \ac{IBM} per l'algoritmo di Shor \cite{ibm_shor_tutorial}. La decisione è motivata da quattro ragioni convergenti.

\textbf{Prima ragione: controllo esplicito del modello.} Qiskit Aer permette di comporre canali depolarizzanti, termici e di readout e di legarli a un gate set dichiarato. Questa flessibilità consente esperimenti controllati e riproducibili; non rende però i preset uniformi della tesi direttamente equivalenti alla calibrazione di una QPU reale.

\textbf{Seconda ragione: copertura delle primitive necessarie.} Qiskit Aer espone nello stesso ambiente i canali depolarizzanti, termici, di readout e coerenti richiesti dal disegno. Altri framework, fra cui Cirq \cite{cirq2024} e PennyLane \cite{pennylane2022}, offrono anch'essi simulazione rumorosa; la scelta non pretende di stabilire una graduatoria generale, ma riduce le conversioni rispetto al codice Qiskit adottato.

\textbf{Terza ragione: disponibilità dell'implementazione di riferimento.} Il circuito dell'algoritmo di Shor utilizzato in questa tesi è tratto da un'implementazione open-source in Qiskit \cite{github_shor}. Questo punto di partenza, validato dalla letteratura \cite{skosana2021}, ha reso naturale mantenere Qiskit come framework unico lungo l'intero pipeline sperimentale.

\textbf{Quarta ragione: allineamento con le fonti implementative.} Diversi lavori usati come riferimento adottano Qiskit \cite{skosana2021, shor_hpec2023, ml_mitigation2023}. Condividere il framework facilita il confronto delle API e dei circuiti, ma non rende direttamente comparabili risultati ottenuti con gate set, versioni, transpiler o noise model differenti.

\begin{table}[H]
\centering
\caption{Confronto tra i principali framework per la simulazione quantistica con rumore.}
\label{tab:framework_confronto}
\begin{tabularx}{\textwidth}{lXXX}
\toprule
\textbf{Framework} & \textbf{Sviluppatore} & \textbf{Simulatore con rumore} & \textbf{Target hardware primario} \\
\midrule
Qiskit + Aer        & IBM Quantum           & Completo (T1, T2, depolarz., readout) & IBM superconduttori \\
Cirq                & Google Quantum AI     & Parziale                               & Google superconduttori \\
PennyLane           & Xanadu                & Limitato                               & Vari (plug-in) \\
Forest / PyQuil     & Rigetti               & Parziale                               & Rigetti superconduttori \\
\bottomrule
\end{tabularx}
\end{table}

% -----------------------------------------------
\subsection{Qiskit Aer: il Simulatore}
\label{sec:qiskit_aer}

Qiskit Aer è il modulo di simulazione classica di Qiskit. Il suo componente principale è \texttt{AerSimulator}, un simulatore ad alte prestazioni che supporta più metodi di simulazione selezionabili a seconda della dimensione del circuito e del tipo di analisi richiesta.

\subsubsection{Metodi di Simulazione}

\texttt{AerSimulator} espone i seguenti metodi principali tramite il parametro \texttt{method}:

\begin{itemize}
    \item \texttt{statevector} — simulazione del vettore di stato completo. Richiede $2^n$ ampiezze complesse in memoria ed è il metodo usato dalla campagna rumorosa principale su $N=15$.

    \item \texttt{density\_matrix} — evolve la matrice densità $\rho$ anziché il vettore di stato, permettendo di modellare canali di rumore non unitari in modo esatto. Richiede $4^n$ elementi e riduce il limite pratico a circa 25 qubit.

    \item \texttt{stabilizer} — simulazione efficiente di circuiti Clifford puri (senza porte non-Clifford come $T$ o $R_z$ arbitrarie). Scala a migliaia di qubit ma non è applicabile all'algoritmo di Shor, che utilizza rotazioni di fase arbitrarie.

    \item \texttt{matrix\_product\_state} — rappresenta lo stato come rete di tensori; senza troncamenti imposti resta una rappresentazione controllata, ma il costo cresce con l'entanglement. È usato in M8/M13 per contenere il costo delle molte repliche.
\end{itemize}

La scelta è quindi per campagna, non globale: \texttt{statevector} per la prima fase N15 e \texttt{matrix\_product\_state} per la sensibilità M8/M13. Le istanze N21/N35 sono validate idealmente ma escluse dalle campagne rumorose generali.

\subsubsection{Modelli di Rumore}

Il modulo \texttt{qiskit\_aer.noise} fornisce le primitive per costruire modelli di rumore compositi. Ogni errore è rappresentato come un canale quantistico che viene applicato ai gate specificati, con la possibilità di assegnare errori diversi a qubit diversi (rumore eterogeneo) o un errore uniforme a tutti i qubit (\textit{all-qubit error}). I tipi di errore disponibili includono:

\begin{itemize}
    \item \texttt{depolarizing\_error} — errore depolarizzante su gate a 1 o 2 qubit, parametrizzato da $\lambda$ nella convenzione Aer; la probabilità Pauli non-identità è $3\lambda/4$ o $15\lambda/16$.
    \item \texttt{thermal\_relaxation\_error} — modella il rilassamento verso lo stato fondamentale ($T_1$) e la perdita di coerenza di fase ($T_2$), con un tempo di gate configurabile.
    \item \texttt{ReadoutError} — matrice di confusione che descrive la probabilità di misurare $\ket{0}$ come $\ket{1}$ e viceversa.
    \item Errori coerenti — applicazione di una rotazione indesiderata su ogni gate, modellabile tramite operatori unitari arbitrari.
\end{itemize}

Il modello di rumore usato nella prima fase combina errori depolarizzanti, rilassamento termico e readout error con parametri uniformi. È un modello illustrativo per isolare i meccanismi, non uno snapshot di \texttt{ibm\_marrakesh}: non include coupling map, routing, errori per qubit, crosstalk o deriva temporale. La configurazione dettagliata è riportata nel Capitolo~\ref{chap:obiettivi}.

\subsubsection{Transpilazione}

Prima dell'esecuzione, la funzione \texttt{transpile} riscrive il circuito nella base esplicita \texttt{rz/sx/x/cx} e applica ottimizzazioni strutturali. La campagna usa il livello di ottimizzazione 2 e il seed del transpiler 20260819; questa coppia, insieme all'hash della sequenza compilata, è registrata in ogni artefatto. La base esplicita impedisce che porte composte restino opache e sfuggano al canale di rumore previsto.

% -----------------------------------------------
\subsection{Accelerazione GPU con Qiskit Aer}
\label{sec:gpu}

La libreria \texttt{qiskit-aer-gpu} estende \texttt{AerSimulator} con il supporto per l'accelerazione tramite \ac{GPU} NVIDIA attraverso CUDA. Il backend \texttt{statevector} è interamente eseguito sulla \ac{GPU}: il vettore di stato viene allocato nella memoria VRAM e le operazioni di gate sono eseguite come moltiplicazioni matrice-vettore parallelizzate su migliaia di core CUDA.

\textbf{Motivazione pratica.} L'accelerazione era stata valutata per la generazione dei dataset e per le campagne replicate. Non vengono però riportati fattori di accelerazione ipotetici come risultati sperimentali, poiché l'ambiente canonico finale usa la \ac{CPU}.

\textbf{Specifiche hardware.} La \ac{GPU} prevista è una NVIDIA GeForce RTX 4070 con 12 GB di memoria VRAM e 5888 core CUDA (architettura Ada Lovelace). La sola allocazione di uno statevector \texttt{complex128} avrebbe un tetto teorico di 29 qubit, inferiore nella pratica per memoria di lavoro e overhead. Questo limite dimensionale non garantisce inoltre la fattibilità: circuiti molto profondi e traiettorie rumorose moltiplicano il costo temporale, ed è proprio ciò che esclude N21/N35 dalla campagna rumorosa.\footnote{La memoria del solo vettore è $2^n \times 16$ byte. Con 12 GB decimali, $\log_2(12 \times 10^9 / 16)\simeq29.48$, quindi il massimo intero teorico è $n=29$.}

\textbf{Esito: l'accelerazione non è stata utilizzabile.} Quanto sopra descrive il piano iniziale. In sede di implementazione il backend \ac{GPU} si è rivelato inutilizzabile nell'ambiente \ac{WSL} adottato --- il driver CUDA riconosce la scheda, ma il backend C++ di Aer solleva un errore a runtime sull'esecuzione dei circuiti --- e \textbf{tutti gli esperimenti riportati in questa tesi sono stati eseguiti su \ac{CPU}}. La diagnosi completa e le conseguenze sui tempi di esecuzione sono nella Sezione~\ref{sec:gpu_limit}. Il vincolo di memoria resta quello indicato, riferito però alla memoria di sistema anziché alla VRAM.

% -----------------------------------------------
\subsection{scikit-learn: la Libreria per il Classificatore}
\label{sec:scikit}

Il Metodo 2 richiede l'addestramento di un classificatore binario. La libreria scelta è scikit-learn \cite{scikit_learn}, una delle librerie di machine learning più consolidate nell'ecosistema Python scientifico.

La motivazione principale è pragmatica: scikit-learn espone un'interfaccia unificata (\texttt{fit} / \texttt{predict}) per un ampio catalogo di classificatori — Random Forest, \ac{SVM}, \ac{MLP}, e altri — permettendo di confrontarli sullo stesso dataset con minime modifiche al codice. Questo facilita la selezione sperimentale del modello migliore senza dover introdurre dipendenze da framework più pesanti (TensorFlow, PyTorch) per un classificatore che, per il problema in esame, si prevede non richieda architetture profonde.

% -----------------------------------------------
\subsection{Ambiente di Esecuzione: WSL con Ubuntu}
\label{sec:wsl}

Le simulazioni vengono eseguite in ambiente Linux tramite \ac{WSL} (versione 2) su Windows 11. La scelta è motivata dalla necessità di mantenere la compatibilità con l'ecosistema Linux di Qiskit e delle librerie scientifiche Python, garantendo al contempo l'integrazione con i driver CUDA per \ac{WSL} forniti da NVIDIA.

\ac{WSL} 2 esegue un kernel Linux completo in una macchina virtuale leggera. L'ambiente Python è isolato in un virtual environment dedicato, garantendo la riproducibilità delle dipendenze e la compatibilità delle versioni tra le librerie usate; le versioni esatte sono versionate nel file \texttt{requirements.txt} della repository. È in questo stesso ambiente che l'accesso alla \ac{GPU} si è rivelato non praticabile (Sezione~\ref{sec:gpu_limit}).

Con l'ambiente così definito, il Capitolo~\ref{chap:metodologia} descrive come Qiskit, Aer e scikit-learn si integrano nell'architettura del pipeline sperimentale: dalla costruzione del circuito di Shor alla generazione del dataset di addestramento, fino all'inferenza del classificatore.

\end{onehalfspace}
